\section{Interference Structure, Complementarity, and Protocol Reconstruction of Delayed-Choice Phenomena}

\subsection{Projection Interference and Formal Correspondence of Moiré Structure}

When cut-and-project maps high-dimensional periodic structure to low dimensions, if the readout kernel has finite width, the statistical structures exhibited in low dimensions generally manifest as quasi-periodic fringes, beat patterns, and moiré-type interference textures. This paper distinguishes ``wave-like'' and ``particle-like'' properties as protocol effects:
\begin{itemize}
  \item When the readout kernel is broad (low resolution), multiple microstate equivalence classes superpose, and output exhibits continuous or fringe-like statistics;
  \item When the readout kernel is narrow (high resolution), equivalence classes are split, and output tends toward discrete event sequences.
\end{itemize}
This distinction does not rely on introducing randomness at the ontological layer, but is induced jointly by the degenerate structure of $\Fold_m$ and readout kernel width.

\subsection{Non-Commutative Structure and Complementarity}

If the scan is generated by shift operator $U$ and phase operator $V$, satisfying the Weyl relation
\[
UV=e^{2\pi i\alpha}VU,
\]
there exists a structural limitation on simultaneous precise determination of scan phase and readout localization. This paper interprets complementarity as: within a single protocol, attempting to simultaneously minimize ``scan localization error'' and ``readout localization error'' is subject to lower bounds imposed by non-commutative structure; this fact can be rewritten as lower bounds on sample complexity of statistical identification tasks or error trade-off relations, thus becoming a testable protocol-layer proposition.

\subsection{Protocol Explanation of Delayed-Choice-Like Phenomena}

Delayed-choice-like experiments can be expressed in this framework as: after readout records are generated, change the projection parameter (such as changing the readout kernel, changing the acceptance window, or changing the selection rule for complementary phase offset), thereby changing how ``the same static ontology decomposes under low-dimensional readout into equivalence classes.'' Under this explanation:
\begin{itemize}
  \item No retrocausal influence on the ontological layer needs to be introduced;
  \item Reconstruction occurs at the protocol layer: what changes is how the readout operator decomposes and classifies the static information structure, not the static structure itself.
\end{itemize}
The testable point is that this reconstruction should correspond to computable equivalence class reallocation and degenerate structure change, and should manifest statistically as predictable reweight migration.
